\documentclass[11pt,a4paper]{article}

\usepackage[margin=1in]{geometry}
\usepackage{booktabs}
\usepackage{graphicx}
\usepackage{hyperref}
\usepackage{amsmath}

\title{Market-Basket Analysis on the IMDB Movies Dataset}
\author{Student name and student ID}
\date{Algorithms for Massive Data, academic year 2025/26}

\begin{document}
\maketitle

\section*{Declaration}

I/We declare that this material, which I/We now submit for assessment, is
entirely my/our own work and has not been taken from the work of others, save
and to the extent that such work has been cited and acknowledged within the
text of my/our work, and including any code produced using generative AI
systems. I/We understand that plagiarism, collusion, and copying are grave and
serious offences in the university and accept the penalties that would be
imposed should I engage in plagiarism, collusion or copying. This assignment,
or any part of it, has not been previously submitted by me/us or any other
person for assessment on this or any other course of study.

\section{Introduction}

Describe the goal of the project: finding frequent actor co-appearances in the
IMDB Movies Dataset using a market-basket analysis technique studied in the
course.

\section{Dataset And Scope}

State that the dataset is the IMDB Movies Dataset from Kaggle. Describe the
columns used, the number of movies considered, and whether the full dataset or
a selected view was used.

\section{Data Organization}

Explain how movies were represented as baskets and how actor names from
\texttt{Star1}, \texttt{Star2}, \texttt{Star3}, and \texttt{Star4} were
represented as items.

\section{Preprocessing}

Describe missing-value handling, duplicate handling within a basket, text
normalization, and any filtering choices.

\section{Algorithms And Implementation}

Describe the selected frequent-itemset technique. Explain the algorithmic idea,
the implementation choices, and any parameters such as minimum support.

\section{Scalability}

Discuss how the solution scales with the number of baskets, number of distinct
items, maximum itemset size, and support threshold. Include empirical runtime
or candidate-count measurements.

\section{Experiments}

Describe the experimental setup, thresholds, hardware or Colab runtime, and the
outputs collected.

\section{Results And Discussion}

Present the main frequent itemsets and interpret them. Discuss whether the
patterns are meaningful and how this method relates to practical affinity
analysis in domains such as product analytics or financial services.

\section{Limitations}

Discuss limitations of the dataset, small basket size, actor-credit ordering,
and the fact that the IMDB setting is only an analogy for market-basket
transactions.

\section{Conclusion}

Summarize what the implementation shows and what could be extended.

\bibliographystyle{plain}
\bibliography{references}

\end{document}

